% Part: computability
% Chapter: machines-computations
% Section: halting-problem

\documentclass[../../../include/open-logic-section]{subfiles}

\begin{document}

\olfileid{tur}{und}{hal}
\olsection{د درېدنې مسئله}

\begin{explain}
فرض کړئ موږ د ټيورينګ ماشين د بيانونو کومه شمېره ثابته کړې ده۔ نو
هر ټيورينګ ماشين يو \emph{شاخص} اخلي: د ټيورينګ ماشين د بيانونو په
شمېره~$M_1$، $M_2$، $M_3$، \dots{} کښې د هغه ځاے۔

موږ پوهېږو چې د ټيورينګ په وسيله نۀ محاسبه کېدونکې تابعې بايد موجودې
وي: د ټيورينګ ماشين د بيانونو سټ---او له دې سره د ټيورينګ ماشينونو
سټ---!!{enumerable} دے، خو له~$\Nat$ څخه~$\Nat$ ته د ټولو تابعو سټ
داسې نۀ دے۔ خو موږ د نامحاسبه کېدونکو تابعو ټاکلې بېلګې هم موندلے
شو۔ يوه داسې تابع د درېدنې تابع ده۔
\end{explain}

\begin{defn}[د درېدنې تابع]
 د \emph{درېدنې تابع}~$h$ داسې تعريفېږي:
\[
h(e,n) =
\begin{cases}
  \text{0} & \text{که ماشين~$M_e$ د ننوت $n$ دپاره ونۀ درېږي} \\
  \text{1} & \text{که ماشين~$M_e$ د ننوت $n$ دپاره ودرېږي}
\end{cases}
\]
\end{defn}

\begin{defn}[د درېدنې مسئله]
د \emph{درېدنې مسئله} د دې ټاکلو مسئله ده چې (د هر~$e$ او~$n$ دپاره)
ايا ټيورينګ ماشين~$M_e$ د~$n$ دانو کرښو پر ننوت درېږي که نه۔
\end{defn}

\begin{explain}
موږ د دې په ښودلو ثابتوو چې~$h$ د ټيورينګ په وسيله محاسبه کېدونکې
نۀ ده چې ور پورې تړلې تابع~$s$ د ټيورينګ په وسيله محاسبه کېدونکې
نۀ ده۔ دا ثبوت پر دې حقيقت تکيه کوي چې هر هغه څه چې يو ټيورينګ ماشين
ئې محاسبه کولے شي د يوه منضبط ټيورينګ ماشين په وسيله هم محاسبه کېدے
شي (\olref[mac][dis]{sec})، او پر دې حقيقت چې دوه ټيورينګ ماشينونه
سره نښلول کېدے شي، څو يو واحد ماشين جوړ کړي (\olref[mac][cmb]{sec})۔
\end{explain}

\begin{defn} تابع~$s$ داسې تعريفېږي:
\[
s(e) =
\begin{cases}
  \text{0} & \text{که ماشين~$M_e$ د ننوت $e$ دپاره ونۀ درېږي} \\
  \text{1} & \text{که ماشين~$M_e$ د ننوت $e$ دپاره ودرېږي}
\end{cases}
\]
\end{defn}

\begin{lem}
تابع~$s$ د ټيورينګ په وسيله محاسبه کېدونکې نۀ ده۔
\end{lem}

\begin{proof}
د تناقض دپاره فرضوو چې تابع~$s$ د ټيورينګ په وسيله محاسبه کېدونکې
ده۔ بيا به داسې ټيورينګ ماشين~$S$ موجود وے چې~$s$ محاسبه کوي۔ د
عموميت له زيان پرته فرضولے شو چې کله~$S$ درېږي، د لومړۍ مربع د
لوستلو پر مهال درېږي (يعنې منضبط دے)۔ دا ماشين له بل ماشين~$J$ سره
``نښلېدے'' شي چې که پر ننوت~$0$ پيل شي درېږي (يعنې که په پيلني
حالت کښې د پټې د پاي-نښې ښي لور ته مربع لوستلو پر مهال~$\TMblank$
ولولي)، او که نۀ وي ښي لور ته روانېږي او هېڅکله نۀ درېږي۔
$S \concat J$، يعنې هغه ماشين چې د~$S$ له~$J$ سره په نښلولو جوړېږي،
يو ټيورينګ ماشين دے، نو د کوم~$e$ دپاره~$M_e$ دے (يعنې په شمېره کښې
يو ځاے راځي)۔ $M_e$ د~$e$ دانو~$\TMstroke$ پر ننوت پيل کړئ۔ دوه
امکانونه دي: يا~$M_e$ درېږي يا نۀ درېږي۔
\begin{enumerate}
\item فرض کړئ~$M_e$ د~$e$ دانو~$\TMstroke$ پر ننوت درېږي۔ بيا~$s(e)
  = 1$۔ نو کله چې~$S$ پر~$e$ پيل شي، د وت په توګه پر پټه له يوه
  $\TMstroke$ سره درېږي۔ بيا~$J$ پر داسې پټه پيلېږي چې يو~$\TMstroke$
  لري۔ په هغه حالت کښې~$J$ نۀ درېږي۔ خو~$M_e$ هماغه ماشين~$S
  \concat J$ دے، نو بايد کټ مټ هغه څه وکړي چې لومړے~$S$ او ورپسې~$J$
  ئې کوي (يعنې په دې حالت کښې ښي لور ته روان شي او هېڅکله ونۀ درېږي)۔
  نو~$M_e$ د~$e$ دانو~$\TMstroke$ پر ننوت نۀ شي درېدې۔

\item اوس فرض کړئ~$M_e$ د~$e$ دانو~$\TMstroke$ پر ننوت نۀ درېږي۔
  بيا~$s(e) = 0$، او کله چې~$S$ پر ننوت~$e$ پيل شي، له تشې پټې سره
  درېږي۔ کله چې~$J$ پر تشه پټه پيل شي، سملاسي درېږي۔ بيا هم~$M_e$
  هغه څه کوي چې لومړے~$S$ او ورپسې~$J$ ئې کوي، نو~$M_e$ بايد د~$e$
  دانو~$\TMstroke$ پر ننوت ودرېږي۔
\end{enumerate}
په هر حالت کښې له خپل فرض سره تناقض ته رسېږو۔ دا ښيي چې داسې ټيورينګ
ماشين~$S$ نشي موجودېدے: $s$ د ټيورينګ په وسيله محاسبه کېدونکې نۀ
ده۔
\end{proof}

\begin{thm}[د درېدنې د مسئلې ناحلي]
\ollabel{thm:halting-problem} د درېدنې مسئله ناحل ده، يعنې تابع~$h$
د ټيورينګ په وسيله محاسبه کېدونکې نۀ ده۔
\end{thm}

\begin{proof}
فرض کړئ~$h$ د ټيورينګ په وسيله محاسبه کېدونکې وے، مثلاً د ټيورينګ
ماشين~$H$ په وسيله۔ موږ له~$H$ څخه داسې ټيورينګ ماشين جوړولے شو چې
$s$ محاسبه کړي: لومړے د ننوت يوه کاپي جوړه کړئ (چې په~$\TMblank$
نښه بېلېږي)۔ بيا بېرته پيل ته لاړ شئ او~$H$ وچلوئ۔ موږ په څرګند ډول
داسې ماشين جوړولے شو چې لومړے کار کوي
(\cref{tur:mac:dis:prob:copier} وګورئ)، او که~$H$ موجود وے، له داسې
کاپي‌کوونکي ماشين سره مو هغه ``نښلولے'' او داسې نوے ماشين مو ترې
جوړولے شو چې ټاکي ايا~$M_e$ پر ننوت~$e$ درېږي؛ يعنې~$s$ محاسبه کوي۔
خو موږ لا ښودلې ده چې داسې هېڅ ماشين نشي موجودېدے۔ نو~$h$ هم د
ټيورينګ په وسيله محاسبه کېدونکې نۀ ده۔
\end{proof}

\begin{prob}
د درې-درېدنې مسئله د داسې پرېکړيزې کړنلارې د ورکولو مسئله ده چې
وټاکي ايا په خپلسر غوره شوے ټيورينګ ماشين پر داسې پټه چې له دې پرته
تشه وي، د
درې دانو~$\TMstroke$ پر ننوت درېږي که نه۔ ثابته کړئ چې د درې-درېدنې
مسئله ناحل ده۔
\end{prob}

\begin{prob}
وښيئ چې که د ټيورينګ ماشين او ننوت د هغو جوړو~$M_e$ او~$n$ دپاره
چې~$e \neq n$ وي د درېدنې مسئله حل کېدونکې وي، نو د هغو حالتونو
دپاره هم حل کېدونکې ده چې~$e = n$ وي۔
\end{prob}

\begin{prob}
موږ ثابته کړه چې که ننوت داسې عدد~$e$ وي چې د ټولو ټيورينګ ماشينونو
د يوې شمېرې له لارې ټيورينګ ماشين~$M_e$ پېژني، نو د درېدنې مسئله
ناحل ده۔ که پر ځاے ئې د~\olref[tur][und][enu]{sec} د ټيورينګ
ماشينونو بيانونه نېغ په نېغه د ننوت په توګه ومنو، نو څه؟ ايا داسې
ټيورينګ ماشين موجودېدے شي چې د درېدنې مسئله پرېکړه کړي خو د شاخصونو
پر ځاے د ټيورينګ ماشينونو بيانونه د ننوت په توګه واخلي؟ څرګنده کړئ
چې ولې يا ولې نه۔
\end{prob}

\begin{prob} وښيئ چې \emph{قسمي} تابع~$s'$ چې داسې تعريفېږي:
  \[
  s'(e) =
  \begin{cases}
    \text{1} & \text{که ماشين~$M_e$ د ننوت $e$ دپاره ودرېږي}\\
    \text{ناتعريفه} & \text{که ماشين~$M_e$ د ننوت $e$ دپاره ونۀ درېږي}
  \end{cases}
  \]
  د ټيورينګ په وسيله محاسبه کېدونکې \emph{ده}۔
\end{prob}

\end{document}
